\section{Experimental setup}
\label{sec:setup}

\subsection{Datasets}
\Cref{tab:datasets} summarises the three datasets. Each is replayed in timestamp order. From each we take 400{,}000
flows evenly spaced over time, deterministically, so that every method and seed sees the same data. The first
30\% form the training period. The last 20\% of the training period is the reference set (\cref{sec:problem}),
and the remaining 280{,}000 flows form 70 windows of 4{,}000 flows. Rows without a valid timestamp or label are
skipped when reading.

\textbf{CIC-IDS2017}~\cite{sharafaldin2018} covers five working days of a testbed network. We use the
\texttt{TrafficLabelling} CSVs, the variant that keeps timestamps; their afternoon times are written without
AM/PM, and we restore them from the capture schedule. The training period (Monday and Tuesday morning) contains
only benign traffic and FTP-Patator, so every attack in the stream is new.

\textbf{CSE-CIC-IDS2018} covers ten days with a larger network and attack set. We use the ``Processed Traffic
Data for ML'' CSVs, thinned to one row in eight per file before sampling to keep the replay tractable. Seven
attack classes occur in the training period and seven appear later. Several day files of the published CSVs are
truncated at 1{,}048{,}575 rows, a spreadsheet limit, which we leave as is.

\textbf{UNSW-NB15}~\cite{moustafa2015unsw} was captured on two days, about four weeks apart, with a different
tool chain. Every class already occurs on the first day, so no class would be new in a chronological replay. As is
common in open-set NIDS work, we remove four attack classes (DoS, Reconnaissance, Shellcode and Worms) from the
training period. They appear in the stream as unknown attacks. The other classes are known, but their February
traffic differs from the January training traffic.

\begin{table}[t]
\centering
\caption{Datasets (after replay sampling: 400k flows, 70 windows of 4{,}000 flows).}
\label{tab:datasets}
\small
\begin{tabular}{@{}lccc@{}}
\toprule
 & CIC-IDS2017 & CSE-CIC-IDS2018 & UNSW-NB15 \\
\midrule
Source flows & 2.83\,M & 16.2\,M & 2.54\,M \\
Days covered & 5 & 10 & 2 \\
Classes (incl.\ benign) & 15 & 15 & 10 \\
Attack classes in training & 1 & 7 & 5 \\
New in stream & 12 & 7 & 4 (held out) \\
Model input features & 79 & 78 & 97 \\
\bottomrule
\end{tabular}
\end{table}

\textbf{Preprocessing.} Identifier columns (flow ID, source and destination address, source port, timestamps;
for UNSW-NB15 also the destination port) are dropped, so the model cannot memorise hosts. Non-finite values are replaced by the training-period median.
Numeric features are transformed with $\operatorname{sign}(v)\log(1+|v|)$ to tame heavy tails, then standardised
with training-period statistics. Categorical fields (UNSW-NB15's protocol, state and service) are one-hot encoded
with an ``unseen'' bucket. Nothing is fitted on the stream.

\subsection{Label regime and classifier}
The main budget is $B = 10\%$ of each window (400 labels), and the low budget $B = 1\%$ (40 labels). The delay is
$\delta = 1$ window, and \cref{sec:res-sens} varies it. Every method that uses labels receives exactly the same
number; INSOMNIA uses none by design.

All methods use the same classifier, scikit-learn's histogram gradient boosting with 150 iterations and L2
regularisation 1.0. Classes with fewer than five labelled flows are left out of training and stay \unk{}. All
methods also use the same open-world detector (\cref{sec:unknown}), except in the detector ablation. The
differences between methods therefore come from the adaptation strategy alone. \Cref{sec:res-sens} repeats the
main comparison with a two-layer MLP.

\subsection{Baselines}
\label{sec:baselines}
We compare with four generic strategies and four methods re-implemented from the literature. All replay memories
hold 20{,}000 flows and are initialised from the training period.
\begin{description}
  \item[Static.] The initial model, never updated: the lower bound.
  \item[Periodic retrain.] From scratch every five windows, on the training data plus the last five labelled
        windows; always deployed.
  \item[Continual replay.] A uniform reservoir memory; retrained whenever labels arrive; always deployed.
  \item[Guarded selective.] Retrained only when a drift detector (Kolmogorov--Smirnov test, population stability
        index or mean-confidence drop) fires or the \unk{} rate exceeds 15\%. The newest labelled window is used as
        a hold-out, the candidate must pass the same gain and forgetting checks as in \nova{}, and a deployed model
        can be rolled back. It shows what a guard achieves without \nova{}'s other components.
  \item[Uncertainty sampling.] Least-confidence active learning on top of continual replay. It is the standard
        query baseline of the active-learning works above and INSOMNIA's ``uncertainty sampling + oracle'' upper
        bound~\cite{andresini2021insomnia}.
  \item[INSOMNIA]~\cite{andresini2021insomnia}, multi-class: after each window, the least-confident 50\% of its
        flows are labelled by a nearest-centroid estimator and added to a pseudo-labelled memory; no analyst
        labels are used.
  \item[NetGuard PS]~\cite{gupta2025netguard}: NetGuard's density-aware prior selection, with one diagonal
        Gaussian mixture (up to eight components) on the labelled memory and one on the new window. The flows with
        the largest log-density ratio are labelled. The generative minority augmentation, which needs GPU hours
        per cycle, is not reproduced.
  \item[CBRS replay]~\cite{chrysakis2020online,amalapuram2023augmented}: the class-balancing reservoir memory used
        for NIDS replay, with uniformly sampled labels, retrained whenever labels arrive.
\end{description}

\subsection{Development and held-out data}
All design decisions, including the components of \nova{}, its hyper-parameters and the CIC-IDS2017 ablation,
were made on CIC-IDS2017. CSE-CIC-IDS2018 and UNSW-NB15 were not used for any design decision. We only checked
their class composition with one run of the static model, which revealed that no UNSW-NB15 class is new without a
hold-out. Both datasets are evaluated with the configuration frozen on CIC-IDS2017. The analyses that we ran on
them after the main results (the ablations and the low budget) change nothing in the method; we mark them as such.

\subsection{Experiments, statistics and implementation}
\Cref{tab:experiments} maps the experiments to the research questions. Each configuration runs with three seeds;
the seed changes model initialisation, label sampling and clustering, not the data. We report the mean and
standard deviation over seeds. To compare \nova{} with another method, we take the per-window difference of their
seed-averaged OW-F1 (or false-positive rate) and report its mean with a 95\% moving-block bootstrap
interval~\cite{kunsch1989jackknife}: blocks of five consecutive windows, 10{,}000 resamples. Consecutive windows are
autocorrelated, and resampling blocks rather than single windows respects this; a test that treats windows as
independent, such as a paired Wilcoxon test, would overstate certainty. We call a difference significant only if
the interval excludes zero.

Everything is implemented in Python (scikit-learn, NumPy) inside an open-source research console that stores
the configuration, data hash, model versions, adaptation decisions and per-flow predictions of every run. Tree
models run on a 16-core CPU; the MLP uses PyTorch 2.14.1 with CUDA 13.0 on an RTX 3060 GPU with deterministic
kernels. In total, the article reports \TBD{N} runs.

\begin{table}[t]
\centering
\caption{Experiments. Unless stated otherwise: budget 10\%, delay 1, gradient boosting, three seeds.}
\label{tab:experiments}
\small
\resizebox{\textwidth}{!}{%
\begin{tabular}{@{}llll@{}}
\toprule
RQ & Experiment & Datasets & Section \\
\midrule
RQ1 & random vs.\ chronological split, static classifier & all & \ref{sec:res-split} \\
RQ2 & nine methods & all & \ref{sec:res-main} \\
RQ3 & budget 1\%, seven label-using methods & all & \ref{sec:res-budget} \\
RQ4 & six one-component ablations; centroid vs.\ $k$-NN detector & all; CIC-IDS2017 & \ref{sec:res-ablation} \\
RQ5 & delay 0 and 3; stricter and more lenient thresholds; MLP & CIC-IDS2017 & \ref{sec:res-sens} \\
-- & per-class results; representation-learning probe & all; CIC-IDS2017 & \ref{sec:res-perclass}, \ref{sec:res-meta} \\
\bottomrule
\end{tabular}}
\end{table}
